% paper2_grf_essay.tex
% Gravity as Information Recovery: The Refractive Index is the Petz Fidelity
% Gravity Research Foundation Essay (2026)
% Author: Sheng-Kai Huang
% Compile: pdflatex paper2_grf_essay.tex (run 3x)

\documentclass[prl,twocolumn,superscriptaddress,floatfix,nofootinbib]{revtex4-2}

\usepackage{amsmath}
\usepackage{amssymb}
\usepackage{hyperref}
\usepackage{booktabs}

% Macros
\newcommand{\kB}{k_{\mathrm{B}}}
\newcommand{\Tr}{\mathrm{Tr}}
\newcommand{\Petz}{\mathcal{R}}
\newcommand{\chan}{\mathcal{N}}
\newcommand{\tauasym}{\tau}
\newcommand{\ceff}{c_{\mathrm{eff}}}
\newcommand{\rs}{r_{\mathrm{s}}}
\newcommand{\Sgrav}{\Sigma_{\mathrm{grav}}}
\newcommand{\Fbound}{F_{\mathrm{bound}}}
\newcommand{\nref}{n_{\mathrm{grav}}}

\begin{document}

\title{Gravity as Information Recovery:\\
The Gravitational Refractive Index is the Petz Fidelity Bound}

\author{Sheng-Kai Huang}
\email{akai@fawstudio.com}
\affiliation{Independent Researcher}

\date{\today}

\begin{abstract}
Three quantities from three fields turn out to be the same number.
From general relativity: the coordinate speed of light
$\ceff/c = e^{\Phi/c^2}$.
From thermodynamics: the entropy production cost
$e^{-\Sgrav/2}$ with $\Sgrav = 2|\Phi|/c^2$.
From quantum information: the Petz recovery fidelity bound
$\Fbound \geq e^{-\Sigma/2}$.
The identification $\ceff/c = e^{-\Sgrav/2} = \Fbound$
states that light slowing in a gravitational field \emph{is}
information loss, and the gravitational refractive index \emph{is}
the inverse recovery fidelity.
Three independent routes---modular flow, gravitational
Landauer principle, and quantum channel analysis---converge on
$\Sgrav = -\ln(-g_{00})$.
If the Petz bound is saturated, the metric is uniquely
the exponential metric $g_{00} = -e^{-\rs/r}$, which has no
event horizon ($\tauasym < 1$ everywhere) and agrees with
Schwarzschild through second post-Newtonian order.
Strong-field deviations---$4.6\%$ larger shadow, $\sim 4\%$
shifted quasinormal modes, and gravitational-wave
echoes at $\Delta t \approx 4\,\rs/c$---are within reach of
current observatories.
\end{abstract}

\maketitle

%=======================================================================
\section{Introduction}
%=======================================================================

In 1911, Einstein observed that the speed of light varies in a
gravitational field~\cite{Einstein1911}.
Dicke later showed that general relativity (GR) can be exactly
reformulated as flat-spacetime optics with a gravitational
refractive index~\cite{Dicke1957}:
in a static field with Newtonian potential $\Phi < 0$,
\begin{equation}
  \frac{\ceff}{c} = e^{\Phi/c^2}\,.
  \label{eq:ceff}
\end{equation}
Meanwhile, quantum information theory has established that
any quantum channel $\chan$ with entropy production
$\Sigma \equiv D(\rho\|\sigma) - D(\chan(\rho)\|\chan(\sigma))$
admits a universal recovery bound~\cite{Petz1988,JRSWW2018}:
\begin{equation}
  F\!\bigl(\rho,\;\Petz_\sigma \circ \chan(\rho)\bigr)
  \;\geq\; e^{-\Sigma/2}\,,
  \label{eq:petz_bound}
\end{equation}
where $\Petz_\sigma$ is the (rotated) Petz recovery
map~\cite{Petz1988} and $F$ is the Uhlmann fidelity.
The Petz map is the \emph{unique} retrodiction functor
satisfying Bayesian consistency~\cite{Parzygnat2023}---the
quantum generalization of Bayes' rule.
In a companion paper~\cite{Huang2026}, we showed that the temporal
asymmetry parameter $\tauasym = 1 - F \in [0,1]$ provides an
operational measure of the arrow of time, with
$\tauasym = 0$ for closed systems and
$\tauasym \leq 1 - e^{-\Sigma/2}$ in general.

Here we connect these two results.
Equations~\eqref{eq:ceff} and~\eqref{eq:petz_bound} share the same
exponential structure. This is not a coincidence: they are the same
physical statement in different languages.

%=======================================================================
\section{The Triple Identification}
\label{sec:identification}
%=======================================================================

We define the dimensionless gravitational entropy
production as
\begin{equation}
  \Sgrav \;\equiv\; \frac{2|\Phi|}{c^2}
  \;=\; \frac{\rs}{r}\,,
  \label{eq:Sgrav}
\end{equation}
where $\rs = 2GM/c^2$ is the Schwarzschild radius.
This is an \emph{informational} entropy production---the
cost of retrodiction from infinity to radius $r$---not a
Clausius entropy. The Clausius entropy production for a
photon in Tolman thermal equilibrium vanishes, since
$T(r)\sqrt{-g_{00}} = T_\infty$~\cite{Tolman1930}. The
quantity $\Sgrav$ persists even in equilibrium: it measures
the informational cost of climbing the gravitational
potential.

Three independent arguments converge on
Eq.~\eqref{eq:Sgrav}, all sharing the mathematical core
$\Sgrav = -\ln(-g_{00})$:

\emph{(a) Modular flow} [known].---In algebraic QFT on
Schwarzschild, the quantum relative entropy (QRE)
between the vacuum and a coherent excitation is
computable via modular theory~\cite{DorauMuch2025}.
The Tolman blueshift and energy redshift together give
a fractional QRE loss
$(D_{\mathrm{in}} - D_{\mathrm{out}})/D_{\mathrm{in}}
= \rs/r$, exactly.

\emph{(b) Gravitational Landauer principle} [known].---The
Tolman-modified cost of erasing one bit at radius $r$
exceeds the cost at infinity by
$1/\sqrt{-g_{00}}$~\cite{Herrera2020}.
The dimensionless excess erasure cost yields
$\Sgrav = -\ln(-g_{00})$, which equals $\rs/r$ for the
exponential metric and
$-\ln(1 - \rs/r) \approx \rs/r + O((\rs/r)^2)$ for
Schwarzschild.

\emph{(c) Quantum channel} [known].---Modeling gravitational
redshift as a pure-loss bosonic channel with intensity
transmissivity $\eta = -g_{00}$, the QRE decrease is
$\Sigma = -\ln\eta = \rs/r$~\cite{AhmadiFuentes2014}.

Substituting Eq.~\eqref{eq:Sgrav} into the Petz
bound~\eqref{eq:petz_bound} gives $\Fbound = e^{-\Sgrav/2}
= e^{\Phi/c^2}$. Comparing with Eq.~\eqref{eq:ceff}:
\begin{equation}
  \boxed{\;\frac{\ceff}{c}
  \;=\; e^{-\Sgrav/2}
  \;=\; \Fbound\,.\;}
  \label{eq:triple}
\end{equation}
This is the central result [new synthesis]: the coordinate
speed of light \emph{is} the Petz recovery fidelity bound.
Equivalently, the gravitational refractive index is the
inverse fidelity:
$\nref = c/\ceff = 1/\Fbound = e^{+\Sgrav/2}$.

The physical picture is transparent.
A photon climbing out of a gravitational potential loses
information---its quantum state degrades due to
interaction with the gravitational channel.
The recovery fidelity $\Fbound$ quantifies
how well the original state can be retrodicted.
At the would-be event horizon, $\Fbound \to 0$ and
$\ceff \to 0$: complete information loss
coincides with the vanishing of the coordinate speed
of light. The event horizon \emph{is} an information
horizon, literally.

Table~\ref{tab:numerical} verifies the identification
across four orders of magnitude in gravitational
compactness.
The agreement is exact for the exponential metric and
holds through $O(\rs/r)$ for Schwarzschild; the first
discrepancy is at $O((\rs/r)^3)$, which is
$\sim 10^{-18}$ in the solar
system~\cite{Bertotti2003}---far below measurability.

\begin{table}[b]
\caption{\label{tab:numerical}%
Numerical verification of the triple identification.
$\Phi/c^2$ is the dimensionless potential,
$\Sgrav = 2|\Phi|/c^2$, and
$\Fbound = e^{-\Sgrav/2} = \ceff/c$.}
\begin{ruledtabular}
\begin{tabular}{lcccc}
  \textbf{Location} & $\Phi/c^2$ & $\Sgrav$ & $\Fbound$ & $\ceff/c$ \\
  \colrule
  Earth
    & $-7{\times}10^{-10}$
    & $1.4{\times}10^{-9}$ & $1{-}7{\times}10^{-10}$
    & $1{-}7{\times}10^{-10}$ \\
  Sun
    & $-2.1{\times}10^{-6}$
    & $4.2{\times}10^{-6}$ & $1{-}2.1{\times}10^{-6}$
    & $1{-}2.1{\times}10^{-6}$ \\
  NS
    & $-0.2$
    & $0.4$ & $0.819$ & $0.819$ \\
  Horizon
    & $-\infty$
    & $\infty$ & $0$ & $0$ \\
\end{tabular}
\end{ruledtabular}
\end{table}

%=======================================================================
\section{Consequences}
\label{sec:consequences}
%=======================================================================

\emph{Model-independent no-horizon theorem.}---The Petz
bound alone, before specifying $\Sgrav$, has a
consequence: if gravitational redshift is a quantum
channel with \emph{finite} entropy production, then
$F \geq e^{-\Sigma/2} > 0$, so $\tauasym < 1$.
An event horizon requires $F = 0$
(complete information loss), which demands
$\Sigma \to \infty$---impossible for a finite-mass
object at finite distance.
The Petz bound forbids event horizons
model-independently~\cite{JRSWW2018};
the specific form of $\Sgrav$ is not needed.

\emph{The exponential metric.}---If the Petz bound
is \emph{saturated}---$\sqrt{-g_{00}} = e^{-\Sgrav/2}$
exactly---the metric is uniquely
determined~\cite{Papapetrou1954,Yilmaz1958}:
\begin{equation}
  ds^2 = -e^{-\rs/r}\, c^2 dt^2
       + e^{+\rs/r}\!\left(dr^2 + r^2 d\Omega^2\right),
  \label{eq:exp_metric}
\end{equation}
in isotropic coordinates.
Since $e^{-\rs/r} > 0$ for all $r > 0$, there is no
event horizon: $\tauasym < 1$ everywhere.
This metric has a long independent pedigree:
Papapetrou~\cite{Papapetrou1954},
Yilmaz~\cite{Yilmaz1958},
Rosen's bimetric theory~\cite{Rosen1973},
and three scalar-field families within GR that all
reduce to it when scalar charge equals gravitational
mass~\cite{Makukov2020}.
Boonserm et~al.\ classified it as a traversable
wormhole with throat at
$R_{\text{thr}} = (e/2)\,\rs$~\cite{Boonserm2018}.
It shares PPN parameters $\gamma = \beta = 1$ with
Schwarzschild and passes all solar-system
tests~\cite{Bertotti2003}.
Our contribution is the information-theoretic motivation:
the exponential metric is the geometry that \emph{saturates}
the fundamental bound on quantum information recovery.

\emph{Caveat on saturation.}---No known channel exactly
saturates the JRSWW bound~\cite{JRSWW2018} for
$\Sigma > 0$; approximate saturation holds in weak fields
with gap $O((\rs/r)^2)$.
Whether nature selects saturation (exponential metric) or
strict inequality (a metric between exponential and
Schwarzschild) is empirical.
The exponential metric is \emph{not} a vacuum solution of
the Einstein equations; it requires a phantom scalar
source~\cite{Makukov2020,Boonserm2018}.
We motivate the \emph{metric} from information theory,
not field equations---an important distinction from
the Yilmaz program~\cite{Yilmaz1958},
whose modified field equations face well-documented
difficulties~\cite{Misner1995}.

\emph{Gravitational decoherence.}---Pikovski
et~al.~\cite{Pikovski2015} showed that gravitational time
dilation universally dephases composite quantum systems in
spatial superposition.
In the $\tauasym$ framework~\cite{Huang2026}, the resulting
temporal asymmetry is $\tauasym_{\mathrm{grav}}(t)
= 1 - e^{-N\Gamma t^2}$, where $\Gamma \propto
(\Delta E)^2(\Delta\Phi)^2/(\hbar^2 c^4)$.
The dephasing rate is controlled by
$\Delta\Phi/c^2 \sim \Sgrav/2$: the gravitational entropy
production governs both the static (refractive index)
and dynamic (decoherence) aspects of gravitational
information loss.
Bonder et~al.~\cite{Bonder2015} noted that the Pikovski
effect is frame-dependent; our framework is consistent
with this, since $\tauasym$ depends on the observer's
channel, not on an absolute decoherence rate.

%=======================================================================
\section{Relation to Prior Work}
\label{sec:prior}
%=======================================================================

The logic here is complementary to the Jacobson
program~\cite{Jacobson1995}.
Jacobson derives \emph{field equations} from
thermodynamic entropy applied to local Rindler horizons:
$\delta Q = T\delta S$ $\to$ Einstein equations $\to$
Schwarzschild.
We derive the \emph{metric component} directly from
information recoverability:
$\Fbound = e^{-\Sgrav/2}$ $\to$ $g_{00}$ $\to$
exponential metric.
Both programs share the premise that geometry is
information-theoretic; they differ in what is derived
first (field equations vs.\ metric) and consequently
arrive at different strong-field predictions.
The Dorau--Much result~\cite{DorauMuch2025}---deriving
the semiclassical Einstein equations from quantum
relative entropy via modular theory---suggests that the
two programs share a common informational foundation and
may ultimately converge.

In holographic settings, the Petz map reconstructs bulk
quantum fields from boundary data~\cite{Cotler2019}, with
fidelity $F_{\text{bulk}} \geq e^{-\Delta S_{\text{gen}}/2}$,
where $\Delta S_{\text{gen}}$ is the change in generalized
entropy.
Our identification~\eqref{eq:triple} extends this
structure beyond AdS/CFT: the bulk metric
\emph{itself} encodes the recovery fidelity, with
$\sqrt{-g_{00}}$ playing the role of $\Fbound$.
The Crooks fluctuation theorem has recently been
established in curved spacetimes~\cite{BassoCeleri2025},
confirming that the Crooks--Petz connection of
Paper~I~\cite{Huang2026} extends to general-relativistic
settings---entropy production in curved spacetime is
observer-dependent and quantified by $\Sgrav$.

%=======================================================================
\section{Testable Predictions}
\label{sec:predictions}
%=======================================================================

The exponential and Schwarzschild metrics agree through
second post-Newtonian order but diverge in strong
fields, yielding three testable predictions ordered by
near-term feasibility.

\emph{(1) Shadow size.}---The exponential photon sphere
at $r_{\text{ph}} = \rs$ (isotropic coordinates)
yields a critical impact parameter
$b_{\text{crit}} = 2e\, m \approx 5.44\, m$
versus Schwarzschild's $3\sqrt{3}\,m \approx 5.20\,m$
($m = GM/c^2$): a $4.6\%$ larger shadow.
For M87*, this gives $41.5\;\mu$as versus
$39.7\;\mu$as; the EHT measurement of
$42 \pm 3\;\mu$as~\cite{EHT2019} is consistent with
both but closer to the exponential value.
Next-generation EHT, targeting sub-$\mu$as precision,
can distinguish the two.

\emph{(2) Quasinormal modes.}---In the eikonal limit,
$f_{\text{QNM}}^{\text{exp}}/f_{\text{QNM}}^{\text{Schw}}
= 3\sqrt{3}/(2e) \approx 0.956$,
a $\sim 4.4\%$ downward shift.
This estimate should be treated with caution: the
exponential photon sphere coincides with the wormhole
throat, where the effective potential differs
qualitatively from Schwarzschild.
WKB calculations for a generalized exponential
wormhole find a $-7.8\%$ shift~\cite{NathSarma2024}.
Current LIGO ringdown analyses constrain QNM deviations
at the $\sim 5$--$10\%$ level; next-generation
detectors (Einstein Telescope, Cosmic Explorer) will
reach $\sim 1\%$ precision.

\emph{(3) Gravitational-wave echoes.}---The most
distinctive signature.
The exponential metric's traversable wormhole structure
produces echoes---repeated, damped pulses following the
initial ringdown---with time delay
\begin{equation}
  \Delta t_{\text{echo}} \approx \frac{4.2\,\rs}{c}\,,
  \label{eq:echo}
\end{equation}
giving $\sim 2.5$~ms for a $60\,M_\odot$
remnant at $\sim 400$~Hz, squarely in the LIGO band.
This is $\sim 90\times$ shorter than Planck-wall
models~\cite{Cardoso2016}, because the exponential
metric has no horizon and its tortoise coordinate
$r^* = \int e^{\rs/r}\,dr$ remains bounded.
Current echo searches use $\sim 200$~ms
templates; dedicated short-delay searches are needed
to test this prediction.

%=======================================================================
\section{Scope and Limitations}
\label{sec:scope}
%=======================================================================

The identification~\eqref{eq:triple} \emph{can}:
unify quantum error correction, thermodynamic
irreversibility, and gravitational physics under a
single parameter $\tauasym$; provide universal recovery
bounds that hold independently of the specific protocol;
connect the Hayden--Preskill decoder---which \emph{is}
the Petz map~\cite{Cotler2019}---to gravitational
physics; and provide a precise, computable definition
of the gravitational arrow of time: $\tauasym > 0$
whenever $\Sgrav > 0$, with $\tauasym = 0$ if and only
if spacetime is flat.

It \emph{cannot}: explain dark matter ($\Sgrav \sim 10^{-6}$
at galactic scales, far too small);
explain the cosmological constant or dark energy;
derive the full metric tensor $g_{ab}$ from $\tauasym$
(the relation $g_{00} \to \Sgrav \to \tauasym$
is one-directional; $\tauasym$ constrains only $g_{00}$,
not the spatial components);
or extend to dynamical spacetimes without
modification (the covariant formulation
$\Fbound(A \!\to\! B) = |\xi|_A/|\xi|_B$, where
$\xi^a$ is a timelike Killing vector, requires
stationarity).

%=======================================================================
\section{Conclusion}
%=======================================================================

The coordinate speed of light, the informational entropy
production cost, and the Petz recovery fidelity bound
are the same number:
$\ceff/c = e^{-\Sgrav/2} = \Fbound$.
This is a new connection, not a new theory---each piece
is known; the synthesis is new.
If the Petz bound is saturated, the exponential metric
follows uniquely, with testable strong-field deviations
from Schwarzschild.
If saturation fails, the no-horizon theorem persists
model-independently: for any finite $\Sgrav$,
$\tauasym < 1$ and complete information loss cannot occur.

The deepest implication is that the metric tensor does
not merely \emph{resemble} an information-theoretic
quantity---$\sqrt{-g_{00}}$ \emph{is} the recovery
fidelity bound.
Whether nature saturates this bound is a question that
the strong-field observational program of the coming
decade can answer.

\begin{acknowledgments}
The author thanks M.~M.~Wilde for feedback on a
companion manuscript.
\end{acknowledgments}

%=======================================================================
\begin{thebibliography}{25}

\bibitem{Einstein1911}
A.~Einstein,
Ann.\ Phys.\ (Leipzig) \textbf{35}, 898 (1911).

\bibitem{Dicke1957}
R.~H.~Dicke,
Rev.\ Mod.\ Phys.\ \textbf{29}, 363 (1957).

\bibitem{Petz1988}
D.~Petz,
Q.\ J.\ Math.\ \textbf{39}, 97 (1988).

\bibitem{JRSWW2018}
M.~Junge, R.~Renner, D.~Sutter, M.~M.~Wilde, and A.~Winter,
Ann.\ Henri Poincar\'{e} \textbf{19}, 2955 (2018).

\bibitem{Parzygnat2023}
A.~J.~Parzygnat and F.~Buscemi,
Quantum \textbf{7}, 1013 (2023).

\bibitem{Huang2026}
S.-K.~Huang,
``The Arrow of Time from Petz Recovery,''
companion paper (2026).

\bibitem{Tolman1930}
R.~C.~Tolman,
Phys.\ Rev.\ \textbf{35}, 904 (1930).

\bibitem{DorauMuch2025}
L.~Dorau and A.~Much,
Phys.\ Rev.\ Lett.\ (2025);
arXiv:2510.24491.

\bibitem{Herrera2020}
L.~Herrera,
Entropy \textbf{22}, 340 (2020).

\bibitem{AhmadiFuentes2014}
M.~Ahmadi, D.~E.~Bruschi, C.~Sabin, G.~Adesso,
and I.~Fuentes,
Sci.\ Rep.\ \textbf{4}, 4996 (2014).

\bibitem{Bertotti2003}
B.~Bertotti, L.~Iess, and P.~Tortora,
Nature \textbf{425}, 374 (2003).

\bibitem{Papapetrou1954}
A.~Papapetrou,
Z.\ Phys.\ \textbf{139}, 518 (1954).

\bibitem{Yilmaz1958}
H.~Yilmaz,
Phys.\ Rev.\ \textbf{111}, 1417 (1958).

\bibitem{Rosen1973}
N.~Rosen,
Gen.\ Relativ.\ Gravit.\ \textbf{4}, 435 (1973).

\bibitem{Makukov2020}
M.~A.~Makukov and E.~G.~Mychelkin,
Found.\ Phys.\ \textbf{50}, 1346 (2020).

\bibitem{Boonserm2018}
P.~Boonserm, T.~Ngampitipan, A.~Simpson, and M.~Visser,
Phys.\ Rev.\ D \textbf{98}, 084048 (2018).

\bibitem{Misner1995}
C.~W.~Misner,
arXiv:gr-qc/9504050 (1995).

\bibitem{Pikovski2015}
I.~Pikovski, M.~Zych, F.~Costa, and \v{C}.~Brukner,
Nat.\ Phys.\ \textbf{11}, 668 (2015).

\bibitem{Bonder2015}
Y.~Bonder, E.~Okon, and D.~Sudarsky,
Phys.\ Rev.\ D \textbf{92}, 124050 (2015).

\bibitem{Jacobson1995}
T.~Jacobson,
Phys.\ Rev.\ Lett.\ \textbf{75}, 1260 (1995).

\bibitem{Cotler2019}
J.~Cotler, P.~Hayden, G.~Penington, G.~Salton,
B.~Swingle, and M.~Walter,
Phys.\ Rev.\ X \textbf{9}, 031011 (2019).

\bibitem{BassoCeleri2025}
M.~L.~W.~Basso, J.~Maziero, and L.~C.~Celeri,
Phys.\ Rev.\ Lett.\ \textbf{134}, 050406 (2025).

\bibitem{EHT2019}
Event Horizon Telescope Collaboration,
Astrophys.\ J.\ Lett.\ \textbf{875}, L1 (2019).

\bibitem{NathSarma2024}
S.~Nath and D.~Sarma,
arXiv:2401.03738 (2024).

\bibitem{Cardoso2016}
V.~Cardoso, E.~Franzato, and P.~Pani,
Phys.\ Rev.\ Lett.\ \textbf{116}, 171101 (2016).

\end{thebibliography}

\end{document}
